% =====================================================================
\section{Seminarios prácticos}
% =====================================================================
Esta sección resuelve los nueve problemas de la lista \#3 del tema \cite{gonzalez1}, con su numeración
original (31--39), y las cuestiones del parcial de 2025/26 \cite{parcial} que corresponden a la segunda ley.
Cada problema sigue la estructura de los temas 1 y 2: \emph{enunciado}, \emph{idea}, \emph{resolución} paso a
paso con todo el desarrollo simbólico y, cuando aporta, una \emph{observación} y una figura. Salvo
indicación contraria, $R=\SI{8,314}{J/(mol.K)}$; para el Ar, $\bar C_V=\tfrac32R$ y $\gamma=\tfrac53$; para un
diatómico, $\bar C_V=\tfrac52R$ y $\gamma=\tfrac75$.

\subsection{Cómo plantear un problema del Tema 3}

Los problemas del tema piden una de dos cosas: un \emph{rendimiento} o un \emph{cambio de entropía}. Lo
primero es un balance de energía tramo a tramo (Tema 2); lo segundo, un balance de entropía
(tabla~\ref{tab:metodos}).

\begin{table}[htbp]
\centering
\caption{Mapa de decisión: qué método conviene según lo que pide el enunciado.}
\label{tab:metodos}
\small
\begin{tabularx}{\textwidth}{Xlr}
\toprule
Lo que da / pide el enunciado & Método & Problemas\\
\midrule
Ciclo con varios tramos; hallar $\eta$ & M1. Rendimiento de un ciclo & 31--33, 36, P.4(c)\\
$\Delta S$ del sistema (gas, cuerpo, agua\dots) & M2. Camino reversible auxiliar & 34--37, 39, P.3\\
$\Delta S$ de focos, entorno y universo & M3. Balance de entropía & 34, 35, 37, P.3(b), P.4(b)\\
Trabajo máximo, temperatura final óptima & M4. Reversibilidad, $\Delta S_{\mathrm{univ}}=0$ & 38\\
Sistema no ideal: $S(T,V)$, adiabáticas & M5. Ecuaciones $T\,\dd S$ & 37, 39\\
\bottomrule
\end{tabularx}
\end{table}

\begin{recuadro}{Los cinco métodos}
\begin{enumerate}[label=\textbf{\color{navy}M\arabic*.},leftmargin=2.2em,itemsep=4pt]
\item \emph{Rendimiento de un ciclo.} Completar $T$ (y $p$, $V$) en todos los vértices; tabla de $Q$ y $W$ por
  tramo; $W_{\mathrm{neto}}=-\sum Q_i$; $\eta=|W_{\mathrm{neto}}|/Q_{\mathrm{abs}}$ con $Q_{\mathrm{abs}}=\sum_{Q_i>0}Q_i$.
\item \emph{Camino reversible auxiliar.} $\Delta S_{\mathrm{sis}}$ solo depende de los estados extremos: se calcula
  con \eqref{eq:Sgas} o con la tabla~\ref{tab:dS}, aunque el proceso real sea irreversible.
\item \emph{Balance de entropía.} Cada foco: $\Delta S_F=-Q_{\mathrm{sis}}/T_F$; cada cuerpo finito:
  $C\ln(T_f/T_i)$; fuentes de trabajo: 0. Sumar y comprobar $\Delta S_{\mathrm{univ}}\geq0$.
\item \emph{Trabajo máximo.} Imponer $\Delta S_{\mathrm{univ}}=0$ para fijar el estado final y obtener el trabajo
  de la primera ley.
\item \emph{Ecuaciones $T\,\dd S$.} $\dd S=(\dd U+p\,\dd V)/T$ con la ecuación de estado y la energética;
  comprobar la exactitud e integrar. La adiabática reversible es $\dd S=0$.
\end{enumerate}
\end{recuadro}

\subsection{El balance de entropía paso a paso}\label{sec:receta}

\paso{Paso 0: universo} Separar el universo en sistema, focos térmicos, cuerpos finitos y fuentes de trabajo.

\paso{Paso 1: estados} Completar los estados inicial y final de cada parte ($T$, $p$, $V$). Si hay equilibrio
térmico final, $T_f$ sale de la primera ley del conjunto aislado.

\paso{Paso 2: calores reales} Calcular el $Q$ que intercambia cada parte en el proceso \emph{real}, con los
signos de cada una ($Q_F=-Q_{\mathrm{sis}}$).

\paso{Paso 3: entropía del sistema} Por un camino reversible auxiliar entre los mismos estados (M2). No se usa el
$Q$ real salvo que el proceso sea reversible.

\paso{Paso 4: entropía del entorno} Focos, $Q_F/T_F$; cuerpos finitos, $C\ln(T_f/T_i)$; fuentes de trabajo, 0.

\paso{Paso 5: universo} Sumar. Si sale $\Delta S_{\mathrm{univ}}<0$, hay un error de signo; si sale $0$, el proceso
es reversible; si sale positivo, $S_{\mathrm{gen}}=\Delta S_{\mathrm{univ}}$.

\begin{figure}[htbp]
\centering
\begin{tikzpicture}[font=\footnotesize,
   p/.style={draw=navy,fill=fondo,align=center,rounded corners=1pt,inner sep=3pt,text width=2.55cm,minimum height=1.0cm},
   d/.style={draw=acento,fill=white,diamond,aspect=2.4,align=center,inner sep=1pt,font=\scriptsize},
   f/.style={-{Stealth[length=4pt]},navy}]
  \node[p] (s0) at (0,0) {\textbf{0} Partes del\\ universo};
  \node[p] (s1) at (3.15,0) {\textbf{1} Estados\\ inicial y final};
  \node[p] (s2) at (6.3,0) {\textbf{2} Calores\\ reales $Q_i$};
  \node[d] (dd) at (9.6,0) {¿reversible?};
  \node[p,text width=3.0cm] (sA) at (12.55,0.95) {\textbf{3A} Reversible\\ $\Delta S=\int\dbar Q/T$};
  \node[p,text width=3.0cm] (sB) at (12.55,-0.95) {\textbf{3B} Irreversible\\ camino auxiliar (M2)};
  \node[p] (s4) at (9.6,-2.6) {\textbf{4} Entorno\\ $-Q/T_F$, $C\ln\frac{T_f}{T_i}$};
  \node[p] (s5) at (6.3,-2.6) {\textbf{5} $\Delta S_{\mathrm{univ}}\geq0$\\ $=S_{\mathrm{gen}}$};
  \draw[f] (s0) -- (s1); \draw[f] (s1) -- (s2); \draw[f] (s2) -- (dd);
  \draw[f] (dd.north) |- node[pos=0.25,left=2pt,font=\scriptsize\bfseries,acento] {sí} (sA.west);
  \draw[f] (dd.south) |- node[pos=0.25,left=2pt,font=\scriptsize\bfseries,acento] {no} (sB.west);
  \draw[f] (sA.east) -- ++(0.4,0) |- (s4.east);
  \draw[navy] (sB.east) -- ++(0.4,0);
  \draw[f] (s4) -- (s5);
\end{tikzpicture}
\caption{Esquema del balance de entropía. La única bifurcación es la reversibilidad, que decide si la entropía
del sistema se calcula con el calor real o con un camino auxiliar; el entorno siempre se calcula con el calor
real.}
\label{fig:receta}
\end{figure}

\begin{observacion}[El error más frecuente]
Calcular $\Delta S_{\mathrm{sis}}=Q_{\mathrm{real}}/T$ en un proceso irreversible. En la expansión contra
$p_{\mathrm{ext}}$ fija del problema 35, $Q/T=R/2$, pero $\Delta S_{\mathrm{gas}}=R\ln2$: la diferencia,
$R(\ln2-\frac12)$, es justamente la entropía generada.
\end{observacion}

% ---------------------------------------------------------------------
\subsection{Problema 31. Ciclo isotermo, isóbaro e isócoro}
% ---------------------------------------------------------------------
\begin{problema}{31}{Rendimiento de un ciclo de tres tramos}
Una cantidad de Ar, gas perfecto, está a \SI{1}{bar}. Se comprime de forma isoterma y reversible hasta
\SI{10}{bar} y regresa al estado inicial por un isóbaro reversible seguido de un isócoro reversible. Dibuje el ciclo
en el plano $p$--$V$ y calcule su rendimiento.
\end{problema}

\begin{resolucion}
\idea Método M1: completar los tres vértices, calcular $Q$ y $W$ de cada tramo con las fórmulas del Tema 2 y
dividir el trabajo neto entre el calor absorbido.

\paso{Paso 0: vértices} Estado 1: $(p_1,V_1,T_1)$. La isoterma a $T_1$ hasta $10p_1$ da $V_2=V_1/10$. Para volver
a 1 con un isóbaro y un isócoro, el isóbaro a $10p_1$ debe llegar a $V_3=V_1$, y por la ecuación de estado
$T_3=T_1\,\frac{p_3V_3}{p_1V_1}=10\,T_1$. Todo se expresa en unidades de $nRT_1$.

\paso{Paso 1: tramo $1\to2$, isotermo} $\Delta U=0$ y
\[
  Q_{12}=nRT_1\ln\frac{V_2}{V_1}=-nRT_1\ln10=-2{,}303\,nRT_1,\qquad W_{12}=-Q_{12}=+2{,}303\,nRT_1 .
\]

\paso{Paso 2: tramo $2\to3$, isóbaro a $10p_1$}
\[
  Q_{23}=n\bar C_p(T_3-T_2)=\tfrac52\,nR\cdot9T_1=22{,}5\,nRT_1,\qquad
  W_{23}=-10p_1\Bigl(V_1-\frac{V_1}{10}\Bigr)=-9\,nRT_1 .
\]

\paso{Paso 3: tramo $3\to1$, isócoro}
$W_{31}=0$ y $Q_{31}=n\bar C_V(T_1-T_3)=\tfrac32\,nR\,(-9T_1)=-13{,}5\,nRT_1$.

\paso{Paso 4: balance y rendimiento} $W_{\mathrm{neto}}=2{,}303-9=-6{,}697\,nRT_1$ (y $\sum Q=+6{,}697$,
como exige $\Delta U_{\mathrm{ciclo}}=0$). Solo absorbe calor el isóbaro, $Q_{\mathrm{abs}}=22{,}5\,nRT_1$:
\[
  \eta=\frac{|W_{\mathrm{neto}}|}{Q_{\mathrm{abs}}}=\frac{6{,}697}{22{,}5}=\boxed{0{,}298}.
\]
\end{resolucion}

\begin{figure}[htbp]
\centering
\begin{tikzpicture}
\begin{axis}[graf,width=7.4cm,height=5.4cm,xmin=0,xmax=1.2,ymin=0,ymax=11.5,
    xtick={0.1,1},xticklabels={$V_1/10$,$V_1$},ytick={1,10},yticklabels={$p_1$,$10p_1$},
    xlabel={$V$},ylabel={$p$}]
  \fill[fondomedio] plot[domain=1:0.1,samples=60] (\x,{1/\x}) -- (axis cs:1,10) -- cycle;
  \addplot[acento,line width=1.1pt,domain=0.1:1,samples=80] {1/x};
  \addplot[rojo,line width=1.1pt] coordinates {(0.1,10) (1,10)};
  \addplot[navy,line width=1.1pt] coordinates {(1,10) (1,1)};
  \draw[acento,-{Stealth[length=5pt]}] (axis cs:0.40,{1/0.40}) -- (axis cs:0.30,{1/0.30});
  \draw[rojo,-{Stealth[length=5pt]}] (axis cs:0.5,10) -- (axis cs:0.6,10);
  \draw[navy,-{Stealth[length=5pt]}] (axis cs:1,6) -- (axis cs:1,5);
  \fill[black] (axis cs:1,1) circle (1.6pt) node[anchor=north west,font=\scriptsize,inner sep=2pt] {1};
  \fill[black] (axis cs:0.1,10) circle (1.6pt) node[anchor=south east,font=\scriptsize,inner sep=2pt] {2};
  \fill[black] (axis cs:1,10) circle (1.6pt) node[anchor=south west,font=\scriptsize,inner sep=2pt] {3};
  \node[font=\scriptsize,rojo,anchor=south] at (axis cs:0.55,10.2) {$Q_{\mathrm{abs}}$};
  \node[font=\scriptsize,acento,anchor=north east] at (axis cs:0.27,3.2) {isoterma};
  \node[font=\scriptsize,navy] at (axis cs:0.65,6) {$|W|$};
\end{axis}
\end{tikzpicture}
\caption{Problema 31. Compresión isoterma (1--2), isóbaro (2--3) e isócoro (3--1). El ciclo se recorre en sentido
horario: es un motor y el área encerrada es el trabajo neto.}
\label{fig:p31}
\end{figure}

% ---------------------------------------------------------------------
\subsection{Problema 32. Ciclo adiabático, isotermo e isócoro}
% ---------------------------------------------------------------------
\begin{problema}{32}{Rendimiento con una adiabática}
Una cantidad de Ar, gas perfecto, a \SI{1}{bar} se comprime de forma adiabática y reversible hasta \SI{20}{bar} y
regresa al estado inicial por un isotermo reversible seguido de un isócoro reversible. Dibuje el ciclo y calcule su
rendimiento.
\end{problema}

\begin{resolucion}
\idea Método M1 con Poisson en el tramo adiabático.

\paso{Paso 0: vértices} Por Poisson con $(\gamma-1)/\gamma=0{,}4$ y $1/\gamma=0{,}6$:
\[
  \frac{T_2}{T_1}=\Bigl(\frac{p_2}{p_1}\Bigr)^{0,4}=20^{0,4}=3{,}314,\qquad
  \frac{V_2}{V_1}=\Bigl(\frac{p_1}{p_2}\Bigr)^{0,6}=20^{-0,6}=0{,}1657 .
\]
La isoterma a $T_2$ termina en $V_3=V_1$, y el isócoro baja de $T_2$ a $T_1$.

\paso{Paso 1: tramo $1\to2$, adiabático} $Q_{12}=0$ y $W_{12}=n\bar C_V(T_2-T_1)=\tfrac32\cdot2{,}314\,nRT_1=3{,}472\,nRT_1$.

\paso{Paso 2: tramo $2\to3$, isotermo a $T_2$}
\[
  Q_{23}=nRT_2\ln\frac{V_1}{V_2}=3{,}314\,\ln\frac{1}{0{,}1657}\,nRT_1=3{,}314\cdot1{,}797\,nRT_1=5{,}957\,nRT_1=-W_{23}.
\]

\paso{Paso 3: tramo $3\to1$, isócoro} $W_{31}=0$ y $Q_{31}=n\bar C_V(T_1-T_2)=-3{,}472\,nRT_1$.

\paso{Paso 4: rendimiento} $W_{\mathrm{neto}}=3{,}472-5{,}957=-2{,}486\,nRT_1$; solo absorbe la isoterma:
\[
  \eta=\frac{2{,}486}{5{,}957}=\boxed{0{,}418}.
\]
\end{resolucion}

\begin{figure}[htbp]
\centering
\begin{tikzpicture}
\begin{axis}[graf,width=7.4cm,height=5.4cm,xmin=0,xmax=1.2,ymin=0,ymax=22.5,
    xtick={0.1657,1},xticklabels={$V_2$,$V_1$},ytick={1,20},yticklabels={$p_1$,$20p_1$},
    xlabel={$V$},ylabel={$p$}]
  \fill[fondomedio] plot[domain=1:0.1657,samples=60] (\x,{\x^(-5/3)}) -- plot[domain=0.1657:1,samples=60] (\x,{3.314/\x}) -- cycle;
  \addplot[navy,line width=1.1pt,domain=0.1657:1,samples=80] {x^(-5/3)};
  \addplot[rojo,line width=1.1pt,domain=0.1657:1,samples=80] {3.314/x};
  \addplot[acento,line width=1.1pt] coordinates {(1,3.314) (1,1)};
  \draw[navy,-{Stealth[length=5pt]}] (axis cs:0.45,{0.45^(-5/3)}) -- (axis cs:0.36,{0.36^(-5/3)});
  \draw[rojo,-{Stealth[length=5pt]}] (axis cs:0.32,{3.314/0.32}) -- (axis cs:0.42,{3.314/0.42});
  \fill[black] (axis cs:1,1) circle (1.6pt) node[anchor=west,font=\scriptsize,inner sep=3pt] {1};
  \fill[black] (axis cs:0.1657,20) circle (1.6pt) node[anchor=south east,font=\scriptsize,inner sep=2pt] {2};
  \fill[black] (axis cs:1,3.314) circle (1.6pt) node[anchor=south west,font=\scriptsize,inner sep=2pt] {3};
  \node[font=\scriptsize,rojo,anchor=south west] at (axis cs:0.45,8.2) {isoterma $T_2$, $Q_{\mathrm{abs}}$};
  \node[font=\scriptsize,navy,anchor=east] at (axis cs:0.40,3.2) {adiabática};
\end{axis}
\end{tikzpicture}
\caption{Problema 32. Compresión adiabática (1--2), expansión isoterma a $T_2$ (2--3) e isócoro (3--1).}
\label{fig:p32}
\end{figure}

% ---------------------------------------------------------------------
\subsection{Problema 33. Ciclo isóbaro, adiabático e isotermo}
% ---------------------------------------------------------------------
\begin{problema}{33}{Rendimiento con razón de presiones}
Un gas perfecto diatómico realiza un ciclo motor reversible formado por una compresión isóbara, una compresión
adiabática y una expansión isoterma hasta el estado inicial. La razón entre la presión máxima y la mínima es
$r_p=p_{\max}/p_{\min}=5$. Dibuje el ciclo y calcule su rendimiento.
\end{problema}

\begin{resolucion}
\idea Método M1. La isoterma cierra el ciclo, luego une el estado 3 (presión máxima) con el 1 (presión mínima) a
la misma temperatura $T_1$.

\paso{Paso 0: vértices} $p_1=p_2=p_{\min}$ (isóbaro), $p_3=5p_1$ y $T_3=T_1$ (isoterma). Por Poisson en $2\to3$,
con $(\gamma-1)/\gamma=2/7$:
\[
  \frac{T_3}{T_2}=\Bigl(\frac{p_3}{p_2}\Bigr)^{2/7}=5^{2/7}=1{,}584\quad\Longrightarrow\quad T_2=\frac{T_1}{1{,}584}=0{,}6314\,T_1 .
\]

\paso{Paso 1: tramo $1\to2$, isóbaro}
$Q_{12}=n\bar C_p(T_2-T_1)=\tfrac72\,nR\,(-0{,}3686\,T_1)=-1{,}290\,nRT_1$ (cede calor).

\paso{Paso 2: tramo $2\to3$, adiabático} $Q_{23}=0$.

\paso{Paso 3: tramo $3\to1$, isotermo} Expansión de $5p_1$ a $p_1$:
$Q_{31}=nRT_1\ln\frac{p_3}{p_1}=nRT_1\ln5=1{,}609\,nRT_1$ (absorbe).

\paso{Paso 4: rendimiento} $|W_{\mathrm{neto}}|=\sum Q=1{,}609-1{,}290=0{,}319\,nRT_1$, y
\[
  \eta=\frac{0{,}319}{1{,}609}=1-\frac{\frac72(1-5^{-2/7})}{\ln5}=\boxed{0{,}198}.
\]
\end{resolucion}

\begin{figure}[htbp]
\centering
\begin{tikzpicture}
\begin{axis}[graf,width=7.4cm,height=5.4cm,xmin=0,xmax=1.2,ymin=0,ymax=5.8,
    xtick={0.2,0.6314,1},xticklabels={$V_3$,$V_2$,$V_1$},ytick={1,5},yticklabels={$p_{\min}$,$p_{\max}$},
    xlabel={$V$},ylabel={$p$}]
  \fill[fondomedio] (axis cs:1,1) -- (axis cs:0.6314,1) -- plot[domain=0.6314:0.2,samples=60] (\x,{(0.6314/\x)^1.4}) -- plot[domain=0.2:1,samples=60] (\x,{1/\x}) -- cycle;
  \addplot[acento,line width=1.1pt] coordinates {(1,1) (0.6314,1)};
  \addplot[navy,line width=1.1pt,domain=0.2:0.6314,samples=80] {(0.6314/x)^1.4};
  \addplot[rojo,line width=1.1pt,domain=0.2:1,samples=80] {1/x};
  \draw[acento,-{Stealth[length=5pt]}] (axis cs:0.86,1) -- (axis cs:0.76,1);
  \draw[rojo,-{Stealth[length=5pt]}] (axis cs:0.45,{1/0.45}) -- (axis cs:0.55,{1/0.55});
  \fill[black] (axis cs:1,1) circle (1.6pt) node[anchor=north west,font=\scriptsize,inner sep=2pt] {1};
  \fill[black] (axis cs:0.6314,1) circle (1.6pt) node[anchor=north,font=\scriptsize,inner sep=3pt] {2};
  \fill[black] (axis cs:0.2,5) circle (1.6pt) node[anchor=south east,font=\scriptsize,inner sep=2pt] {3};
  \node[font=\scriptsize,rojo,anchor=south west] at (axis cs:0.5,2.2) {isoterma $T_1$, $Q_{\mathrm{abs}}$};
  \node[font=\scriptsize,navy,anchor=east] at (axis cs:0.44,1.3) {adiabática};
\end{axis}
\end{tikzpicture}
\caption{Problema 33. Compresión isóbara (1--2), compresión adiabática (2--3) y expansión isoterma (3--1). El área
encerrada, pequeña frente al calor absorbido, explica el rendimiento bajo.}
\label{fig:p33}
\end{figure}

% ---------------------------------------------------------------------
\subsection{Problema 34. Resistencia sumergida en agua}
% ---------------------------------------------------------------------
\begin{problema}{34}{Disipación eléctrica en un recipiente adiabático}
En un recipiente adiabático con $M_a=\SI{100}{g}$ de agua ($c_a=\SI{4,186}{kJ/(kg.K)}$) a $T_i=\SI{300}{K}$ se sumerge una
resistencia de $M_r=\SI{5}{g}$ ($c_r=\SI{0,85}{kJ/(kg.K)}$) y $R=\SI{50}{\ohm}$. Pasa una corriente $I=\SI{1}{A}$ durante
$t=\SI{5}{min}$ y después se deja que la resistencia alcance el equilibrio con el agua. Determine (a) la temperatura
final; (b) los cambios de energía interna de la resistencia y del agua; (c) los cambios de entropía de la
resistencia, del agua y del universo.
\end{problema}

\begin{resolucion}
\idea La energía eléctrica es trabajo disipativo sobre la resistencia; el conjunto agua más resistencia es
adiabático, así que su $\Delta U$ total es ese trabajo (Tema 2). Las entropías salen de $C\ln(T_f/T_i)$ (M2) y el
universo es la suma (M3).

\paso{Paso 0: datos} $W_{\mathrm{el}}=I^2Rt=1^2\cdot50\cdot300=\SI{15000}{J}$. Capacidades:
$C_a=M_ac_a=\SI{418,6}{J/K}$, $C_r=M_rc_r=\SI{4,25}{J/K}$.

\paso{(a) Temperatura final} Conjunto adiabático y rígido: $\Delta U_a+\Delta U_r=W_{\mathrm{el}}$, con los dos
cuerpos a la misma $T_f$ al final:
\[
  (C_a+C_r)(T_f-T_i)=W_{\mathrm{el}}\quad\Longrightarrow\quad
  T_f=300+\frac{15000}{422{,}85}\ \mathrm K=\boxed{\SI{335,47}{K}}.
\]

\paso{(b) Energías internas} $\Delta U_r=C_r\Delta T=4{,}25\cdot35{,}47=\boxed{\SI{150,8}{J}}$ y
$\Delta U_a=C_a\Delta T=\boxed{\SI{14849}{J}}$. La resistencia recibe los \SI{15000}{J} como trabajo y cede
\SI{14849}{J} de calor al agua.

\paso{(c) Entropías} Con $\ln(T_f/T_i)=\ln1{,}11824=0{,}11176$:
\[
  \Delta S_r=C_r\ln\frac{T_f}{T_i}=\boxed{\SI{0,475}{J/K}},\qquad
  \Delta S_a=C_a\ln\frac{T_f}{T_i}=\boxed{\SI{46,78}{J/K}}.
\]
La fuente eléctrica es una fuente de trabajo y no cambia de entropía, y el recipiente es adiabático:
\[
  \Delta S_{\mathrm{univ}}=\Delta S_r+\Delta S_a=\boxed{\SI{47,26}{J/K}}>0 .
\]
\end{resolucion}

\begin{observacion}[Origen de la entropía]
No entra calor desde fuera, así que toda la entropía es \emph{generada}: la disipación del trabajo eléctrico
y el paso de calor de la resistencia, más caliente durante el proceso, al agua.
\end{observacion}

% ---------------------------------------------------------------------
\subsection{Problema 35. Mismos estados, reversible e irreversible}
% ---------------------------------------------------------------------
\begin{problema}{35}{Expansión isoterma y contra presión fija}
Un mol de gas ideal pasa de $(p_i,V_i)$ a $(p_i/2,2V_i)$ por dos procesos: (a) isotermo reversible y (b)
irreversible contra una presión constante $p_i/2$, en contacto con un foco a $T_i$. Determine en cada caso $Q$, $W$,
$\Delta U$ y $\Delta S$ del sistema y del foco.
\end{problema}

\begin{resolucion}
\idea Los estados son los mismos y están a $T_i$ (porque $p_iV_i=\tfrac{p_i}{2}\,2V_i$): $\Delta U$ y $\Delta S_{\mathrm{gas}}$
son iguales en los dos casos (M2). Cambian $Q$ y $W$, y con ellos la entropía del foco (M3).

\paso{Paso 0: comunes} $\Delta U=0$ (gas ideal, $T$ igual) y, por la tabla~\ref{tab:dS},
$\Delta S_{\mathrm{gas}}=R\ln\frac{2V_i}{V_i}=R\ln2=\SI{5,763}{J/K}$ en los dos casos.

\paso{(a) Reversible}
$W_a=-RT_i\ln2$ y $Q_a=-W_a=RT_i\ln2$. El foco cede ese calor a su temperatura:
\[
  \Delta S_F=-\frac{Q_a}{T_i}=-R\ln2,\qquad \Delta S_{\mathrm{univ}}=R\ln2-R\ln2=\boxed{0}.
\]

\paso{(b) Contra $p_i/2$} Trabajo con la presión exterior (Tema 2, M2):
\[
  W_b=-\frac{p_i}{2}\,(2V_i-V_i)=-\frac{p_iV_i}{2}=-\frac{RT_i}{2},\qquad Q_b=-W_b=\frac{RT_i}{2}.
\]
\[
  \Delta S_F=-\frac{Q_b}{T_i}=-\frac R2,\qquad
  \Delta S_{\mathrm{univ}}=R\Bigl(\ln2-\frac12\Bigr)=8{,}314\cdot0{,}193\ \mathrm{J/K}=\boxed{\SI{1,61}{J/K}}>0 .
\]
\end{resolucion}

\begin{table}[H]
\centering
\caption{Problema 35. Los mismos estados extremos con dos procesos distintos.}
\small
\renewcommand{\arraystretch}{1.5}
\begin{tabular}{lcccccc}
\toprule
Proceso & $W$ & $Q$ & $\Delta U$ & $\Delta S_{\mathrm{gas}}$ & $\Delta S_{\mathrm{foco}}$ & $\Delta S_{\mathrm{univ}}$\\
\midrule
(a) reversible & $-RT_i\ln2$ & $RT_i\ln2$ & 0 & $R\ln2$ & $-R\ln2$ & $0$\\
(b) contra $p_i/2$ & $-RT_i/2$ & $RT_i/2$ & 0 & $R\ln2$ & $-R/2$ & $R(\ln2-\tfrac12)$\\
\bottomrule
\end{tabular}
\end{table}

% ---------------------------------------------------------------------
\subsection{Problema 36. Ciclo con una adiabática irreversible}
% ---------------------------------------------------------------------
\begin{problema}{36}{Temperaturas y entropía generada}
Un mol de gas perfecto diatómico está a $p_1=\SI{1}{atm}$ y $T_1=\SI{300}{K}$. Se comprime reversible y
adiabáticamente hasta $5p_1$ (estado 2); después, en un isóbaro reversible, absorbe \SI{10}{kJ} (estado 3); desde ahí
se expande adiabática e irreversiblemente contra una presión exterior constante $p_1=p_4$ (estado 4); por último,
un isóbaro reversible a $p_1$ cierra el ciclo. (a) Esquematice el ciclo y calcule la temperatura de cada estado.
(b) Calcule la variación de entropía en $3\to4$.
\end{problema}

\begin{resolucion}
\idea $T_2$ sale de Poisson (adiabática reversible); $T_3$, del calor del isóbaro; $T_4$, del balance de energía de la
expansión irreversible contra $p_1$ (Tema 2), no de Poisson. $\Delta S_{34}$ se calcula por un camino reversible
(M2).

\paso{(a) Paso 1: estado 2} Con $(\gamma-1)/\gamma=2/7$:
$T_2=T_1\,5^{2/7}=300\cdot1{,}5838\ \mathrm K=\boxed{\SI{475,1}{K}}$.

\paso{(a) Paso 2: estado 3} Isóbaro: $Q_{23}=n\bar C_p(T_3-T_2)$ con $\bar C_p=\tfrac72R=\SI{29,10}{J/(mol.K)}$:
\[
  T_3=T_2+\frac{Q_{23}}{n\bar C_p}=475{,}1+\frac{10\,000}{29{,}10}\ \mathrm K=475{,}1+343{,}7=\boxed{\SI{818,8}{K}}.
\]

\paso{(a) Paso 3: estado 4} Adiabático ($Q=0$) contra $p_{\mathrm{ext}}=p_1$, y el gas termina en $p_4=p_1$. Balance:
$\Delta U=W$, es decir $n\bar C_V(T_4-T_3)=-p_1(V_4-V_3)$, con $V_4=nRT_4/p_1$ y $V_3=nRT_3/(5p_1)$:
\[
  \tfrac52R\,(T_4-T_3)=-RT_4+\frac{RT_3}{5}
  \quad\Longrightarrow\quad
  \tfrac72T_4=\Bigl(\tfrac52+\tfrac15\Bigr)T_3
  \quad\Longrightarrow\quad
  T_4=\frac{2{,}7}{3{,}5}\,T_3=\boxed{\SI{631,6}{K}}.
\]

\paso{(a) Paso 4: cierre} El isóbaro $4\to1$ a $p_1$ enfría el gas de \SI{631,6}{K} a \SI{300}{K}. Volúmenes
($V=nRT/p$): $V_1=\SI{24,6}{L}$, $V_2=\SI{7,80}{L}$, $V_3=\SI{13,4}{L}$, $V_4=\SI{51,8}{L}$ (figura~\ref{fig:p36}).

\paso{(b) Entropía en $3\to4$} Camino reversible auxiliar entre $(T_3,5p_1)$ y $(T_4,p_1)$ con \eqref{eq:Sgas}:
\[
  \Delta S_{34}=n\bar C_p\ln\frac{T_4}{T_3}-nR\ln\frac{p_4}{p_3}
  =29{,}10\ln\frac{2{,}7}{3{,}5}+8{,}314\ln5
  =-7{,}55+13{,}38=\boxed{\SI{5,83}{J/K}}.
\]
Es positivo, como debe ser en un proceso adiabático irreversible: $\Delta S=S_{\mathrm{gen}}>0$.
\end{resolucion}

\begin{figure}[htbp]
\centering
\begin{tikzpicture}
\begin{axis}[graf,width=8.4cm,height=5.8cm,xmin=0,xmax=58,ymin=0,ymax=6.6,
    xtick={7.80,13.44,24.62,51.83},xticklabels={$V_2$,$V_3$,$V_1$,$V_4$},ytick={1,5},yticklabels={$p_1$,$5p_1$},
    xlabel={$V$ (L)},ylabel={$p$ (atm)},clip=false]
  \addplot[navy,line width=1.1pt,domain=7.80:24.62,samples=80] {(24.62/x)^1.4};
  \addplot[rojo,line width=1.1pt] coordinates {(7.80,5) (13.44,5)};
  \addplot[gris,line width=1.1pt,densely dashed] coordinates {(13.44,5) (51.83,1)};
  \addplot[acento,line width=1.1pt] coordinates {(51.83,1) (24.62,1)};
  \draw[navy,-{Stealth[length=5pt]}] (axis cs:15,{(24.62/15)^1.4}) -- (axis cs:12,{(24.62/12)^1.4});
  \draw[rojo,-{Stealth[length=5pt]}] (axis cs:9.8,5) -- (axis cs:11.4,5);
  \draw[acento,-{Stealth[length=5pt]}] (axis cs:40,1) -- (axis cs:36,1);
  \fill[black] (axis cs:24.62,1) circle (1.6pt) node[anchor=north,font=\scriptsize,inner sep=3pt] {1};
  \fill[black] (axis cs:7.80,5) circle (1.6pt) node[anchor=east,font=\scriptsize,inner sep=3pt] {2};
  \fill[black] (axis cs:13.44,5) circle (1.6pt) node[anchor=south west,font=\scriptsize,inner sep=2pt] {3};
  \fill[black] (axis cs:51.83,1) circle (1.6pt) node[anchor=north,font=\scriptsize,inner sep=3pt] {4};
  \node[font=\scriptsize,rojo,anchor=south] at (axis cs:10.6,5.45) {$Q_{23}=\SI{10}{kJ}$};
  \node[font=\scriptsize,gris,anchor=south west,align=left] at (axis cs:31,3.2) {adiabática\\irreversible};
\end{axis}
\end{tikzpicture}
\caption{Problema 36. La expansión $3\to4$ es irreversible: no tiene trayectoria en el plano $p$--$V$ y solo se
marcan sus extremos (la línea discontinua es esquemática). Los demás tramos son reversibles.}
\label{fig:p36}
\end{figure}

% ---------------------------------------------------------------------
\subsection{Problema 37. Calentamiento con muchos focos y con uno}
% ---------------------------------------------------------------------
\begin{problema}{37}{Entropía del universo con focos}
Un sistema con $U=nAT$ ($A>0$) se calienta reversiblemente a volumen constante de $T_1$ a $T_2$ en contacto con
infinitos focos de temperaturas entre $T_1$ y $T_2$, y después se enfría a volumen constante hasta $T_1$ con un único
foco a $T_1$. (a) Calcule los cambios de entropía del sistema, de los focos y del universo. (b) Repítalo si el
calentamiento se hubiera hecho con un único foco a $T_2$.
\end{problema}

\begin{resolucion}
\idea $C_V=\dd U/\dd T=nA$ es constante; a volumen fijo $\dbar Q=nA\,\dd T$. El sistema hace un ciclo, así que su
$\Delta S$ total es cero; todo el balance está en los focos (M3).

\paso{Paso 0: sistema} En cada tramo, por la tabla~\ref{tab:dS}, $\Delta S_{\mathrm{sis}}=\pm nA\ln(T_2/T_1)$; en el
conjunto, $\Delta S_{\mathrm{sis}}=0$.

\paso{(a) Calentamiento con infinitos focos} En cada instante el sistema y el foco están a la misma $T$: es
reversible. Cada foco cede $nA\,\dd T$ a su temperatura $T$:
\[
  \Delta S_{\mathrm{focos}}=-\int_{T_1}^{T_2}\frac{nA\,\dd T}{T}=-nA\ln\frac{T_2}{T_1},\qquad \Delta S_{\mathrm{univ}}=0 .
\]

\paso{(a) Enfriamiento con un foco a $T_1$} El foco recibe $nA(T_2-T_1)$ a $T_1$:
\[
  \Delta S_{T_1}=\frac{nA(T_2-T_1)}{T_1},\qquad
  \Delta S_{\mathrm{univ}}=nA\Bigl[\frac{T_2}{T_1}-1-\ln\frac{T_2}{T_1}\Bigr]>0,
\]
positivo porque $x-1>\ln x$ para $x=T_2/T_1>1$. Es también el total del apartado (a).

\paso{(b) Calentamiento con un foco a $T_2$} El foco cede $nA(T_2-T_1)$ a $T_2$:
\[
  \Delta S_{T_2}=-\frac{nA(T_2-T_1)}{T_2},\qquad
  \Delta S_{\mathrm{univ}}^{\mathrm{calent.}}=nA\Bigl[\ln\frac{T_2}{T_1}-\Bigl(1-\frac{T_1}{T_2}\Bigr)\Bigr]>0 .
\]
Sumando el enfriamiento, el $\ln$ se cancela:
\[
  \Delta S_{\mathrm{univ}}^{\mathrm{total}}=nA\Bigl[\frac{T_2}{T_1}-2+\frac{T_1}{T_2}\Bigr]=\boxed{nA\,\frac{(T_2-T_1)^2}{T_1T_2}}>0 .
\]
\end{resolucion}

\begin{observacion}[Lectura]
Con infinitos focos el calentamiento no genera entropía: la diferencia de temperatura con el foco es siempre
infinitesimal. Con un solo foco, cada intercambio a través de una diferencia finita de $T$ genera entropía.
\end{observacion}

% ---------------------------------------------------------------------
\subsection{Problema 38. Trabajo máximo entre dos cuerpos}
% ---------------------------------------------------------------------
\begin{problema}{38}{Motor entre cuerpos finitos}
Dos cuerpos idénticos de capacidades $C_1=C_2=C$, a 400 y \SI{300}{K}, se usan como focos de un motor; están a presión
constante y sin cambios de fase. Determine (a) el trabajo máximo y la temperatura de cada cuerpo cuando se obtiene;
(b) el rendimiento en esas condiciones. (Aplicación: $C=\SI{2}{kJ/K}$, es decir, $c=\SI{2}{kJ/(kg.K)}$ para \SI{1}{kg}
de cada cuerpo.)
\end{problema}

\begin{resolucion}
\idea Método M4: el trabajo es máximo si el conjunto es reversible, $\Delta S_{\mathrm{univ}}=0$. El motor hace
ciclos ($\Delta S_{\mathrm{motor}}=0$) y los dos cuerpos acaban a la misma $T_f$, cuando ya no hay diferencia de
temperatura que aprovechar.

\paso{(a) Paso 1: temperatura final}
\[
  \Delta S_{\mathrm{univ}}=C\ln\frac{T_f}{400}+C\ln\frac{T_f}{300}=C\ln\frac{T_f^2}{400\cdot300}=0
  \quad\Longrightarrow\quad T_f=\sqrt{400\cdot300}=\boxed{\SI{346,4}{K}} .
\]

\paso{(a) Paso 2: trabajo} El cuerpo caliente cede $C(400-T_f)$ y el frío recibe $C(T_f-300)$; la diferencia es el
trabajo:
\[
  |W|_{\max}=C(400-T_f)-C(T_f-300)=C\,(700-2T_f)=2000\cdot7{,}18\ \mathrm J=\boxed{\SI{14,4}{kJ}} .
\]

\paso{(b) Rendimiento} $Q_{\mathrm{abs}}=C(400-T_f)=\SI{107,2}{kJ}$ y
\[
  \eta=\frac{700-2T_f}{400-T_f}=\frac{(\sqrt{400}-\sqrt{300})^2}{\sqrt{400}\,(\sqrt{400}-\sqrt{300})}
      =1-\sqrt{\frac{300}{400}}=\boxed{0{,}134}.
\]
\end{resolucion}

\begin{observacion}[Comparación]
El rendimiento es menor que el de Carnot entre las temperaturas iniciales, $1-300/400=0{,}25$, porque los focos se
acercan en temperatura mientras el motor trabaja. Si los cuerpos se pusieran en contacto directo, $T_f=\SI{350}{K}$,
$W=0$ y $\Delta S_{\mathrm{univ}}=C\ln\frac{49}{48}>0$.
\end{observacion}

% ---------------------------------------------------------------------
\subsection{Problema 39. Radiación térmica}
% ---------------------------------------------------------------------
\begin{problema}{39}{Gas de fotones}
La radiación térmica en un volumen $V$ a temperatura $T$ cumple $U=aVT^4$ y $p=U/(3V)$, con $a$ constante. Calcule
(a) la entropía $S(T,V)$; (b) la ecuación de las adiabáticas. Represente en un diagrama $pV$ el ciclo de Carnot de la
radiación y calcule su rendimiento.
\end{problema}

\begin{resolucion}
\idea Método M5: $\dd S=(\dd U+p\,\dd V)/T$ con las dos ecuaciones dadas; la adiabática es $S=\cte$.

\paso{(a) Paso 1: diferencial} $\dd U=4aVT^3\dd T+aT^4\dd V$ y $p\,\dd V=\tfrac13aT^4\dd V$, luego
\[
  \dd S=\frac{\dd U+p\,\dd V}{T}=4aVT^2\,\dd T+\tfrac43aT^3\,\dd V .
\]

\paso{(a) Paso 2: exactitud} $\partial_V(4aVT^2)=4aT^2=\partial_T\bigl(\tfrac43aT^3\bigr)$. \chk

\paso{(a) Paso 3: integración} Respecto de $T$: $S=\tfrac43aVT^3+g(V)$; derivando respecto de $V$,
$\tfrac43aT^3+g'=\tfrac43aT^3$, así que $g'=0$. Tomando $S\to0$ cuando $T\to0$ (tercera ley):
\[
  \boxed{S=\tfrac43\,aVT^3}.
\]

\paso{(b) Adiabáticas} $S=\cte$ da $VT^3=\cte$. Como $p=\tfrac13aT^4$, $T\propto p^{1/4}$ y
$Vp^{3/4}=\cte$, es decir,
\[
  \boxed{pV^{4/3}=\cte}.
\]

\paso{Ciclo de Carnot} Como $p$ solo depende de $T$, las \emph{isotermas son isóbaras}: rectas horizontales en el plano
$pV$ (figura~\ref{fig:p39}). En la isoterma caliente, de $V_1$ a $V_2$:
$Q_h=\Delta U+p\,\Delta V=aT_h^4\Delta V+\tfrac13aT_h^4\Delta V=\tfrac43aT_h^4(V_2-V_1)$, y análogamente
$Q_c=-\tfrac43aT_c^4(V_3-V_4)$. Las adiabáticas dan $V_2T_h^3=V_3T_c^3$ y $V_1T_h^3=V_4T_c^3$, así que
$V_3-V_4=(T_h/T_c)^3(V_2-V_1)$ y
\[
  \frac{Q_c}{Q_h}=-\frac{T_c^4}{T_h^4}\Bigl(\frac{T_h}{T_c}\Bigr)^3=-\frac{T_c}{T_h}
  \quad\Longrightarrow\quad\boxed{\eta=1-\frac{T_c}{T_h}},
\]
el de Carnot, como exige el teorema para cualquier sustancia.
\end{resolucion}

\begin{figure}[htbp]
\centering
\begin{tikzpicture}
% T_h = 1,25 (p_h = 2,441), T_c = 1 (p_c = 1); V1 = 1, V2 = 2, V3 = 3,906, V4 = 1,953
\begin{axis}[graf,width=7.6cm,height=5.4cm,xmin=0,xmax=4.6,ymin=0,ymax=3.1,
    xtick=\empty,ytick={1,2.441},yticklabels={$p_c$,$p_h$},xlabel={$V$},ylabel={$p$}]
  \fill[fondomedio] (axis cs:1,2.441) -- (axis cs:2,2.441) -- plot[domain=2:3.906,samples=60] (\x,{2.441*(2/\x)^(4/3)})
     -- (axis cs:1.953,1) -- plot[domain=1.953:1,samples=60] (\x,{2.441*(1/\x)^(4/3)}) -- cycle;
  \addplot[rojo,line width=1.1pt] coordinates {(1,2.441) (2,2.441)};
  \addplot[navy,line width=1.1pt,domain=2:3.906,samples=80] {2.441*(2/x)^(4/3)};
  \addplot[acento,line width=1.1pt] coordinates {(3.906,1) (1.953,1)};
  \addplot[navy,line width=1.1pt,domain=1:1.953,samples=80] {2.441*(1/x)^(4/3)};
  \draw[rojo,-{Stealth[length=5pt]}] (axis cs:1.4,2.441) -- (axis cs:1.6,2.441);
  \draw[acento,-{Stealth[length=5pt]}] (axis cs:3.0,1) -- (axis cs:2.8,1);
  \fill[black] (axis cs:1,2.441) circle (1.6pt) node[anchor=south east,font=\scriptsize,inner sep=2pt] {1};
  \fill[black] (axis cs:2,2.441) circle (1.6pt) node[anchor=south west,font=\scriptsize,inner sep=2pt] {2};
  \fill[black] (axis cs:3.906,1) circle (1.6pt) node[anchor=north west,font=\scriptsize,inner sep=2pt] {3};
  \fill[black] (axis cs:1.953,1) circle (1.6pt) node[anchor=north east,font=\scriptsize,inner sep=2pt] {4};
  \node[font=\scriptsize,rojo,anchor=south] at (axis cs:1.5,2.55) {isoterma $T_h$};
  \node[font=\scriptsize,acento,anchor=north] at (axis cs:2.9,0.9) {isoterma $T_c$};
  \node[font=\scriptsize,navy,anchor=west] at (axis cs:3.05,1.6) {$pV^{4/3}=\cte$};
\end{axis}
\end{tikzpicture}
\caption{Problema 39. Ciclo de Carnot de la radiación térmica: las isotermas son horizontales porque $p=\tfrac13aT^4$
solo depende de $T$; las adiabáticas son $pV^{4/3}=\cte$.}
\label{fig:p39}
\end{figure}

% ---------------------------------------------------------------------
\subsection{Parcial 2025/26, cuestión 3}
% ---------------------------------------------------------------------
\begin{problema}{P.3}{Entropía del gas y del foco}
Un mol de gas ideal se expande isoterma y reversiblemente desde $(p_i,V_i)$ hasta $p_i/2$, en contacto con un foco a
$T_i$. (a) ¿Cuál es el cambio de energía y de entropía del gas? (b) ¿Cuál es el cambio de entropía del foco?
\end{problema}

\begin{resolucion}
\idea Es el caso (a) del problema 35: mismos estados $(p_i,V_i)\to(p_i/2,2V_i)$ a $T_i$.

\paso{(a)} $\Delta U=0$ (gas ideal isotermo) y $\Delta S_{\mathrm{gas}}=R\ln\frac{V_f}{V_i}=\boxed{R\ln2}$; el gas absorbe
$Q=RT_i\ln2$.

\paso{(b)} El foco cede ese calor a $T_i$: $\Delta S_F=-\dfrac{RT_i\ln2}{T_i}=\boxed{-R\ln2}$, y
$\Delta S_{\mathrm{univ}}=0$, como corresponde a un proceso reversible.
\end{resolucion}

% ---------------------------------------------------------------------
\subsection{Parcial 2025/26, cuestión 4(b) y 4(c)}
% ---------------------------------------------------------------------
\begin{problema}{P.4}{Entropía del universo y teorema de Carnot}
Un gas perfecto monoatómico se calienta adiabática y reversiblemente de $T_1$ a $T_2$ reduciendo su volumen a la
mitad ($1\to2$); se expande en contacto con un foco a $T_2$ hasta su volumen inicial ($2\to3$); y vuelve al estado 1 en
contacto con un foco a $T_1$ por un proceso irreversible a volumen constante ($3\to1$). (b) ¿Cuál es el cambio de
entropía del universo? (c) ¿Se cumple el teorema de Carnot? (Dato: $\frac{x^n-1}{\ln x^n}>1$ para $x>1$, $n>0$.)
\end{problema}

\begin{resolucion}
\idea El gas hace un ciclo, $\Delta S_{\mathrm{gas}}=0$: todo el cambio está en los dos focos (M3). Para (c) se
compara $\eta=|W|/Q_{\mathrm{abs}}$ (M1) con $\eta_C=1-T_1/T_2$.

\paso{Paso 0: datos} Por Poisson con $\gamma-1=\tfrac23$ y $V_1/V_2=2$: $x\equiv T_2/T_1=2^{2/3}=1{,}587$, de modo que
$\ln2=\tfrac32\ln x$. Calores: $Q_{12}=0$; $Q_{23}=nRT_2\ln2$ (isoterma reversible, absorbe);
$Q_{31}=\tfrac32nR(T_1-T_2)$ (isócoro, cede; cuestión 4(a), Tema 2).

\paso{(b) Entropía del universo} Foco a $T_2$: cede $Q_{23}$, $\Delta S_{T_2}=-nR\ln2=-\tfrac32nR\ln x$. Foco a $T_1$:
recibe $\tfrac32nR(T_2-T_1)$, $\Delta S_{T_1}=\tfrac32nR(x-1)$. Así
\[
  \Delta S_{\mathrm{univ}}=\tfrac32nR\,\bigl[(x-1)-\ln x\bigr]=\tfrac32nR\,(0{,}587-0{,}462)=\boxed{0{,}188\,nR}>0,
\]
positivo por el dato del enunciado (con $n=1$ en la potencia, $x-1>\ln x$): el isócoro con una diferencia finita de
temperatura es irreversible.

\paso{(c) Teorema de Carnot} $|W|=Q_{23}+Q_{31}$ y $Q_{\mathrm{abs}}=Q_{23}$:
\[
  \eta=1-\frac{\tfrac32nR(T_2-T_1)}{nRT_2\ln2}=1-\frac{T_2-T_1}{T_2\ln x}=1-\frac{1-1/x}{\ln x}=1-0{,}801=0{,}199 .
\]
Con $\eta_C=1-T_1/T_2=1-1/x=0{,}370$,
\[
  \eta<\eta_C\iff\frac{1-1/x}{\ln x}>\frac1x\iff\frac{x-1}{\ln x}>1,
\]
que es el dato. \textbf{Sí se cumple}: $0{,}199<0{,}370$.
\end{resolucion}

% ---------------------------------------------------------------------
\subsection{Resumen de resultados}
% ---------------------------------------------------------------------
\begin{table}[H]
\centering
\caption{Resultados de los seminarios del Tema 3.}
\small
\renewcommand{\arraystretch}{1.35}
\begin{tabularx}{\textwidth}{lXc}
\toprule
Problema & Resultado & Método\\
\midrule
31 & $\eta=6{,}697/22{,}5=0{,}298$ & M1\\
32 & $T_2=3{,}314\,T_1$;\enspace $\eta=0{,}418$ & M1\\
33 & $T_2=0{,}631\,T_1$;\enspace $\eta=0{,}198$ & M1\\
34 & $T_f=\SI{335,5}{K}$;\enspace $\Delta U_r=\SI{151}{J}$, $\Delta U_a=\SI{14849}{J}$;\enspace $\Delta S_r=\SI{0,475}{J/K}$, $\Delta S_a=\SI{46,8}{J/K}$, $\Delta S_{\mathrm{univ}}=\SI{47,3}{J/K}$ & M2, M3\\
35 & $\Delta S_{\mathrm{gas}}=R\ln2$;\enspace (a) $\Delta S_F=-R\ln2$, $\Delta S_{\mathrm{univ}}=0$;\enspace (b) $\Delta S_F=-R/2$, $\Delta S_{\mathrm{univ}}=\SI{1,61}{J/K}$ & M2, M3\\
36 & $T_2=\SI{475,1}{K}$, $T_3=\SI{818,8}{K}$, $T_4=\SI{631,6}{K}$;\enspace $\Delta S_{34}=\SI{5,83}{J/K}$ & M1, M2\\
37 & (a) $\Delta S_{\mathrm{univ}}=nA[x-1-\ln x]$;\enspace (b) $nA(T_2-T_1)^2/(T_1T_2)$, con $x=T_2/T_1$ & M3\\
38 & $T_f=\sqrt{T_1T_2}=\SI{346,4}{K}$;\enspace $|W|_{\max}=\SI{14,4}{kJ}$;\enspace $\eta=1-\sqrt{3/4}=0{,}134$ & M4\\
39 & $S=\tfrac43aVT^3$;\enspace $pV^{4/3}=\cte$;\enspace $\eta=1-T_c/T_h$ & M5\\
P.3 & $\Delta U=0$, $\Delta S_{\mathrm{gas}}=R\ln2$;\enspace $\Delta S_F=-R\ln2$ & M2, M3\\
P.4 & $\Delta S_{\mathrm{univ}}=\tfrac32nR(x-1-\ln x)=0{,}188\,nR$;\enspace $\eta=0{,}199<\eta_C=0{,}370$ & M3, M1\\
\bottomrule
\end{tabularx}
\end{table}

\begin{recuadro}{Lista de control antes de entregar}
\begin{itemize}[label={\color{acento}$\blacktriangleright$},leftmargin=1.4em,itemsep=4pt]
\item En un rendimiento, ¿se divide entre el calor \emph{absorbido} y no entre el calor neto?
\item ¿Se ha comprobado $\sum Q_i=-W_{\mathrm{neto}}$ en el ciclo?
\item ¿La entropía del sistema se ha calculado por un camino reversible, y no con $Q_{\mathrm{real}}/T$?
\item ¿El calor del foco lleva el signo cambiado, $Q_F=-Q_{\mathrm{sis}}$, y se divide por la $T$ del foco?
\item ¿Sale $\Delta S_{\mathrm{univ}}\geq0$? Si es negativo, hay un error de signo.
\item ¿Se ha usado Poisson solo en adiabáticas \emph{reversibles}?
\end{itemize}
\end{recuadro}

